\section{Produced-Water Screening Results}
\label{sec:screening-results}

The calibrated surrogate returns finite outputs over the current produced-water screening matrix. The generated matrix spans approximately \SIrange{30.2}{61.4}{\percent} one-stage lithium extraction and \SIrange{2.7}{9.8}{\percent} sodium extraction. The nominal MS-2 clean \Smackover row gives \SI{50.66}{\percent} one-stage lithium extraction and \SI{5.67}{\percent} sodium extraction at $O/A=1.0$, corresponding to $\Dli=1.0267$, $\Dna=0.0601$, and $\Slin=17.09$.

\begin{figure}
    \centering
    \includegraphics[width=0.78\linewidth]{calibrated_li_extraction_surface.png}
    \caption{Calibrated Rezaee surrogate lithium extraction surface over lithium feed concentration and organic-to-aqueous ratio.}
    \label{fig:li-extraction-surface}
\end{figure}

\begin{figure}
    \centering
    \includegraphics[width=0.78\linewidth]{calibrated_selectivity_vs_na_li.png}
    \caption{Lithium-to-sodium selectivity over produced-water sodium loading. The high-\NaIon/\LiIon region is an extrapolation relative to the original Rezaee DOE space.}
    \label{fig:selectivity}
\end{figure}

The current result is most useful as an interface product: it supplies extraction, selectivity, distribution coefficients, raffinate loading, and extract loading in a form suitable for staged contactor and costing studies. It should not yet be read as an experimental guarantee for high-\NaIon/\LiIon \Smackover water.
